\documentclass[11pt,a4paper]{article}
% preamble.tex — estilo común de todos los apuntes Froth.
% Cada apunte hace % preamble.tex — estilo común de todos los apuntes Froth.
% Cada apunte hace % preamble.tex — estilo común de todos los apuntes Froth.
% Cada apunte hace \input{preamble} justo después de \documentclass.
\usepackage[margin=2.4cm]{geometry}
\usepackage{xcolor}
\definecolor{litblue}{RGB}{31,79,121}
\definecolor{litgreen}{RGB}{27,94,32}
\definecolor{litorange}{RGB}{230,81,0}
\usepackage{parskip}
\usepackage{titlesec}
\titleformat{\section}{\Large\bfseries\color{litblue}}{\thesection}{0.6em}{}
\titleformat{\subsection}{\large\bfseries\color{litblue}}{\thesubsection}{0.6em}{}
\usepackage{tcolorbox}
\usepackage{fancyhdr}
\pagestyle{fancy}
\fancyhf{}
\fancyhead[L]{\small\color{litblue}Froth — Apuntes de ML}
\fancyhead[R]{\small\thepage}
\renewcommand{\headrulewidth}{0.4pt}
\usepackage[colorlinks=true,linkcolor=litblue,urlcolor=litblue]{hyperref}

% Cabecera de cada apunte: \cabecera{numero}{titulo}
\newcommand{\cabecera}[2]{%
  \begin{tcolorbox}[colback=litblue,coltext=white,colframe=litblue,arc=2mm]
    {\Large\bfseries Apunte #1 — #2}\\[3pt]
    {\small Proyecto Froth · aprender Machine Learning construyendo}
  \end{tcolorbox}\vspace{0.8em}}

% Cajas temáticas
\newtcolorbox{concepto}[1]{colback=blue!4,colframe=litblue,title={Concepto: #1},arc=1.5mm}
\newtcolorbox{practica}{colback=green!5,colframe=litgreen,title={Pruébalo tú (teclado en mano)},arc=1.5mm}
\newtcolorbox{ojo}{colback=orange!8,colframe=litorange,title={Ojo / error típico},arc=1.5mm}
\newtcolorbox{resumen}{colback=gray!8,colframe=gray!60,title={En una frase},arc=1.5mm}

\newcommand{\codigo}[1]{\texttt{#1}}
 justo después de \documentclass.
\usepackage[margin=2.4cm]{geometry}
\usepackage{xcolor}
\definecolor{litblue}{RGB}{31,79,121}
\definecolor{litgreen}{RGB}{27,94,32}
\definecolor{litorange}{RGB}{230,81,0}
\usepackage{parskip}
\usepackage{titlesec}
\titleformat{\section}{\Large\bfseries\color{litblue}}{\thesection}{0.6em}{}
\titleformat{\subsection}{\large\bfseries\color{litblue}}{\thesubsection}{0.6em}{}
\usepackage{tcolorbox}
\usepackage{fancyhdr}
\pagestyle{fancy}
\fancyhf{}
\fancyhead[L]{\small\color{litblue}Froth — Apuntes de ML}
\fancyhead[R]{\small\thepage}
\renewcommand{\headrulewidth}{0.4pt}
\usepackage[colorlinks=true,linkcolor=litblue,urlcolor=litblue]{hyperref}

% Cabecera de cada apunte: \cabecera{numero}{titulo}
\newcommand{\cabecera}[2]{%
  \begin{tcolorbox}[colback=litblue,coltext=white,colframe=litblue,arc=2mm]
    {\Large\bfseries Apunte #1 — #2}\\[3pt]
    {\small Proyecto Froth · aprender Machine Learning construyendo}
  \end{tcolorbox}\vspace{0.8em}}

% Cajas temáticas
\newtcolorbox{concepto}[1]{colback=blue!4,colframe=litblue,title={Concepto: #1},arc=1.5mm}
\newtcolorbox{practica}{colback=green!5,colframe=litgreen,title={Pruébalo tú (teclado en mano)},arc=1.5mm}
\newtcolorbox{ojo}{colback=orange!8,colframe=litorange,title={Ojo / error típico},arc=1.5mm}
\newtcolorbox{resumen}{colback=gray!8,colframe=gray!60,title={En una frase},arc=1.5mm}

\newcommand{\codigo}[1]{\texttt{#1}}
 justo después de \documentclass.
\usepackage[margin=2.4cm]{geometry}
\usepackage{xcolor}
\definecolor{litblue}{RGB}{31,79,121}
\definecolor{litgreen}{RGB}{27,94,32}
\definecolor{litorange}{RGB}{230,81,0}
\usepackage{parskip}
\usepackage{titlesec}
\titleformat{\section}{\Large\bfseries\color{litblue}}{\thesection}{0.6em}{}
\titleformat{\subsection}{\large\bfseries\color{litblue}}{\thesubsection}{0.6em}{}
\usepackage{tcolorbox}
\usepackage{fancyhdr}
\pagestyle{fancy}
\fancyhf{}
\fancyhead[L]{\small\color{litblue}Froth — Apuntes de ML}
\fancyhead[R]{\small\thepage}
\renewcommand{\headrulewidth}{0.4pt}
\usepackage[colorlinks=true,linkcolor=litblue,urlcolor=litblue]{hyperref}

% Cabecera de cada apunte: \cabecera{numero}{titulo}
\newcommand{\cabecera}[2]{%
  \begin{tcolorbox}[colback=litblue,coltext=white,colframe=litblue,arc=2mm]
    {\Large\bfseries Apunte #1 — #2}\\[3pt]
    {\small Proyecto Froth · aprender Machine Learning construyendo}
  \end{tcolorbox}\vspace{0.8em}}

% Cajas temáticas
\newtcolorbox{concepto}[1]{colback=blue!4,colframe=litblue,title={Concepto: #1},arc=1.5mm}
\newtcolorbox{practica}{colback=green!5,colframe=litgreen,title={Pruébalo tú (teclado en mano)},arc=1.5mm}
\newtcolorbox{ojo}{colback=orange!8,colframe=litorange,title={Ojo / error típico},arc=1.5mm}
\newtcolorbox{resumen}{colback=gray!8,colframe=gray!60,title={En una frase},arc=1.5mm}

\newcommand{\codigo}[1]{\texttt{#1}}

\begin{document}
\cabecera{00}{El mapa del proyecto (léeme primero)}

\section{Qué estamos construyendo y por qué}

\textbf{Froth} es una herramienta que, dado un tema científico, descarga cientos de
papers, los ``entiende'' por su contenido y dibuja un \emph{mapa de subtemas}: burbujas
de papers agrupados por significado, conexiones entre grupos, y zonas vacías (los
\emph{gaps} donde puedes aportar algo nuevo en tu tesis).

El proyecto tiene un doble objetivo, y los dos pesan igual:
\begin{enumerate}
  \item \textbf{Aprender Machine Learning de verdad}, pieza a pieza, construyendo.
  \item \textbf{Tener un producto útil} para literature reviews completos (no el
        clásico ``5 papers que confirman lo que ya pensaba'').
\end{enumerate}

\begin{concepto}{pipeline (tubería)}
Todo el sistema es una \textbf{cadena de 5 pasos} donde la salida de uno es la entrada
del siguiente:
\begin{center}
\codigo{[1] PULL $\rightarrow$ [2] EMBED $\rightarrow$ [3] REDUCE $\rightarrow$ [4] CLUSTER $\rightarrow$ [5] VISUALIZE}
\end{center}
\begin{enumerate}
  \item \textbf{Pull}: bajar papers (título, resumen, autores...) desde bases de datos públicas.
  \item \textbf{Embed}: convertir cada resumen en un \emph{vector} (lista de números que captura su significado).
  \item \textbf{Reduce}: aplastar esos vectores de $\sim$768 dimensiones a 2 para poder dibujarlos.
  \item \textbf{Cluster}: agrupar los puntos cercanos $=$ subtemas descubiertos automáticamente.
  \item \textbf{Visualize}: burbujas, mapa mental, detección de gaps, borrador de review.
\end{enumerate}
Cada paso es, a la vez, un concepto de ML que vas a dominar.
\end{concepto}

\section{El mapa de carpetas (qué es cada cosa y en qué orden mirar)}

Dentro de \codigo{Documents\textbackslash Froth\textbackslash} está todo.
\textbf{Regla de oro:} las carpetas \emph{con número} son tuyas (entra); las \emph{sin
número} son la maquinaria (no tocar). La chuleta completa está en el archivo
\codigo{0\_LEEME\_PRIMERO.md} de la raíz.

\textbf{Tuyas (numeradas):}
\begin{enumerate}
  \item \codigo{1\_Apuntes/} --- \textbf{tu biblioteca} (estos PDFs). Léelos en orden: 00, 01, 02...
  \item \codigo{2\_Datos/} --- los datos: \codigo{raw/} (crudos), \codigo{processed/} (tablas
        limpias), \codigo{embeddings/} (vectores). Se llena solo al correr el pipeline.
  \item \codigo{3\_Resultados/} --- las salidas visibles: mapas, reviews.
  \item \codigo{4\_Notebooks/} --- cuadernos de experimentación (Jupyter), cuando lleguen.
  \item \codigo{5\_Bitacoras/} --- decisiones tomadas y tu diario de aprendizaje.
\end{enumerate}

\textbf{Maquinaria (sin número, no tocar):}
\begin{itemize}
  \item \codigo{froth/} --- \textbf{el código fuente} (el motor). Un archivo por paso del
        pipeline: \codigo{config.py} (ajustes) $\rightarrow$ \codigo{pull.py} $\rightarrow$
        \codigo{embed.py} $\rightarrow$ \codigo{cluster.py} $\rightarrow$ \codigo{graph.py}
        $\rightarrow$ \codigo{gaps.py} $\rightarrow$ \codigo{visualize.py}.
  \item \codigo{.venv/} --- la burbuja con Python y sus librerías (ver Apunte 01).
  \item \codigo{tools/} --- el compilador que convierte los apuntes en PDF.
\end{itemize}

\begin{ojo}
¿Por qué \codigo{froth/} \textbf{no} lleva número, si dijimos que todo debe estar ordenado?
Porque Python importa el código por su nombre (\codigo{from froth import pull}) y un nombre
que empieza por número rompe el lenguaje. El \emph{orden} del pipeline no vive en el nombre
de la carpeta sino dentro del código (\codigo{froth/\_\_init\_\_.py}). Los apuntes y tus
carpetas, en cambio, sí van numerados porque nadie los importa: solo se leen.
\end{ojo}

\section{Las 7 fases, de un vistazo}

\begin{enumerate}
  \item[\textbf{F0}] \textbf{Cimientos} --- entorno Python, Git, estructura. \emph{(casi cerrada)}
  \item[\textbf{F1}] \textbf{Pull} --- bajar 200--500 papers de un topic. \emph{(primera versión ya funciona: 126 papers)}
  \item[\textbf{F2}] \textbf{Embeddings} --- texto $\rightarrow$ vectores; buscador semántico. \emph{(siguiente)}
  \item[\textbf{F3}] \textbf{Clusters} --- las burbujas de subtemas (UMAP + HDBSCAN + BERTopic).
  \item[\textbf{F4}] \textbf{Mapa mental} --- grafo de relaciones entre papers y temas.
  \item[\textbf{F5}] \textbf{Gaps + borrador} --- zonas vacías del mapa y texto estructurado de review.
  \item[\textbf{F6}] \textbf{Deep learning} --- PyTorch, transformers, afinar tu propio modelo.
\end{enumerate}

\section{Cómo usar estos apuntes}

Cada vez que aparezca una habilidad o concepto nuevo en el proyecto, se genera un apunte
PDF numerado que se guarda en \codigo{apuntes/}. Tú solo entras ahí y lees en orden.
Cada apunte tiene cajas de colores: \textcolor{litblue}{\textbf{Concepto}} (la idea),
\textcolor{litgreen}{\textbf{Pruébalo tú}} (comandos para teclear tú mismo) y
\textcolor{litorange}{\textbf{Ojo}} (errores típicos).

\begin{resumen}
Froth es una tubería de 5 pasos que convierte ``un tema'' en ``un mapa de la literatura'';
este apunte es tu plano general: carpetas, fases y orden de lectura.
\end{resumen}

\end{document}
